\documentclass[11pt]{article}

\usepackage[a4paper,margin=30mm]{geometry}
\usepackage{amsmath,amssymb,amsthm,mathtools}
\usepackage{enumitem}
\usepackage{microtype}
\usepackage[hidelinks]{hyperref}
\usepackage[nameinlink,noabbrev]{cleveref}

\newtheorem{theorem}{Theorem}[section]
\newtheorem{proposition}[theorem]{Proposition}
\newtheorem{lemma}[theorem]{Lemma}
\newtheorem{corollary}[theorem]{Corollary}
\theoremstyle{definition}
\newtheorem{definition}[theorem]{Definition}
\theoremstyle{remark}
\newtheorem{remark}[theorem]{Remark}

\newcommand{\K}{\mathbb K}
\newcommand{\T}{\mathbb T}
\newcommand{\NA}{\operatorname{NA}}
\newcommand{\ASE}{\operatorname{ASE}}
\newcommand{\calL}{\mathcal L}
\newcommand{\ext}{\operatorname{ext}}
\newcommand{\co}{\operatorname{co}}
\newcommand{\wstar}{w^*}
\newcommand{\norm}[1]{\lVert #1\rVert}
\newcommand{\Ball}{B}
\newcommand{\Sphere}{S}

\title{Residual Norm Attainment from Isolated Norming Orbits\\
\large A corrected and strengthened partial answer to a question of Jung--Mart\'{\i}n--Rueda Zoca}
\author{Research note}
\date{21 June 2026}

\begin{document}
\maketitle

\begin{abstract}
Let $X$ and $Y\ne\{0\}$ be Banach spaces.  Jung, Mart\'{\i}n and Rueda Zoca asked whether residuality of $\NA(X,Y)$ in $\calL(X,Y)$ forces residuality of $\NA(X,\K)$ in $X^*$.  We audit a proposed proof for finite-boundary ranges.  Its domination argument is correct, but its invocation of the unrestricted Kuratowski--Ulam theorem is not justified for possibly nonseparable Banach spaces.  The conclusion nevertheless survives, because a special cylinder-reflection theorem is valid for completely metrizable spaces.

We then prove a stronger result.  It is enough that $Y$ possess a norming set with one orbit, modulo unimodular scalars, uniformly isolated by a vector of $Y$.  The theorem remains valid in any closed operator subspace containing the relevant rank-one operators.  Consequences include all ranges with property quasi-$\beta$, all finite-boundary ranges, and every finite-dimensional range whose dual ball has only countably many extreme-point orbits.  In the latter class, combining the result with the forward implication in the source paper gives an equivalence between scalar and vector-valued residuality.

A literature search through 21 June 2026 did not locate a proof or counterexample in full generality.  The general question therefore remains open in this note.  The claims below should be regarded as a research-level candidate contribution, not as a substitute for peer review.
\end{abstract}

\section{The question and the literature status}

For Banach spaces $X$ and $Y$ over $\K\in\{\mathbb R,\mathbb C\}$, write
\[
 \NA(X,Y)=\bigl\{T\in\calL(X,Y):\text{there is }x\in\Sphere_X
       \text{ with }\norm{Tx}=\norm{T}\bigr\}.
\]
The question in Remark~4.2 of Jung--Mart\'{\i}n--Rueda Zoca \cite{JMRZ} is:
\begin{quote}
If $\NA(X,Y)$ is residual in $\calL(X,Y)$ for some nonzero Banach space $Y$, must $\NA(X,\K)$ be residual in $X^*$?
\end{quote}
Here ``residual'' means that the complement is meagre.

The source paper gives a positive answer under additional separability assumptions: if $X$ and $Y^*$ are separable, residuality of $\NA(X,Y)$ implies density of $\ASE(X,Y)$ \cite[Theorem~4.1]{JMRZ}; density of $\ASE(X,Y)$ implies density of the strongly exposing functionals on $X$ \cite[Proposition~1.5]{JMRZ}, and hence scalar residuality.  Related nonseparable cases are recorded in \cite[Remark~4.3]{JMRZ}.

The broader question whether residuality of $\NA(X,Y)$ always implies density of $\ASE(X,Y)$ was still explicitly described as open in 2025 \cite[after Corollary~2.18]{RSE}.  Searches by the exact wording of the scalar question, by title citations of \cite{JMRZ}, and through recent work on norm-attaining and range-strongly-exposing operators did not reveal a later full solution or a counterexample.  This is evidence, not a proof of novelty.

\section{Audit of the finite-boundary argument}

The proposed argument has the right main idea.  Let a finite set $\Phi\subset \Sphere_{Y^*}$ norm $Y$ by
\[
 \norm{y}=\max_{\phi\in\Phi}|\phi(y)|,
\]
and suppose that $\phi_0\in\Phi$ and $y_0\in\Sphere_Y$ satisfy
\[
 \phi_0(y_0)=1,
 \qquad
 \max_{\phi\in\Phi\setminus\{\phi_0\}}|\phi(y_0)|<1.
\]
Writing $Z=\ker\phi_0$, every operator has the unique form
\[
 T=y_0\otimes f+H,
 \qquad f=\phi_0\circ T\in X^*,\quad H\in\calL(X,Z).
\]
On a suitable open product cylinder, the $\phi_0$-coordinate has strictly larger norm than all the other coordinates.  One should explicitly insert the resulting equality
\[
 \norm{T}=\max_{\phi\in\Phi}\norm{\phi\circ T}
          =\norm{\phi_0\circ T}=\norm f.
\]
With that equality, the equivalence between norm attainment of $T$ and of $f$ is correct.

There is, however, a gap in the category paragraph as originally written.  The usual Kuratowski--Ulam theorem in the required section form is guaranteed, for example, when the sectioned factor has a countable $\pi$-base.  An arbitrary nonseparable Banach space need not satisfy this hypothesis.  Thus one cannot simply invoke the unrestricted theorem.

Fortunately, only a much more special statement is needed, and that statement is true.

\begin{lemma}[Cylinder reflection]\label{lem:cylinder}
Let $P$ and $Q$ be completely metrizable spaces, with $P\ne\varnothing$.  For every $A\subset Q$,
\[
 A\text{ is meagre in }Q
 \quad\Longleftrightarrow\quad
 P\times A\text{ is meagre in }P\times Q.
\]
Consequently, if $U$ and $V$ are nonempty open subsets of Banach spaces and $A\subset U$, then
\[
 A\times V\text{ residual in }U\times V
 \quad\Longrightarrow\quad
 A\text{ residual in }U.
\]
\end{lemma}

\begin{proof}
The equivalence is the cylinder theorem for completely metrizable spaces; see \cite[Lemma~4.3]{Brian}.  The easy implication follows by taking products of nowhere dense sets with $P$.  The converse is a special product-category fact that remains valid without separability, even though the full section theorem need not be available.  Open subsets of Banach spaces are completely metrizable in compatible metrics, so the final assertion follows by applying the equivalence to the complement of $A$.
\end{proof}

Thus the candidate finite-boundary theorem is correct after replacing its Kuratowski--Ulam paragraph by \cref{lem:cylinder}.  The polyhedral specialization is also correct.  The next section proves a strictly stronger version.

\section{The isolated-norming-orbit theorem}

Let
\[
 \T=\{\lambda\in\K:|\lambda|=1\}.
\]
For $\phi\in Y^*$, write $\T\phi=\{\lambda\phi:\lambda\in\T\}$.

\begin{definition}\label{def:ino}
A norming set $\Gamma\subset\Sphere_{Y^*}$ has an \emph{isolated norming orbit} if there are $\phi_0\in\Gamma$ and $y_0\in\Sphere_Y$ such that
\[
 \phi_0(y_0)=1
 \quad\text{and}\quad
 \beta:=\sup\bigl\{|\gamma(y_0)|:
       \gamma\in\Gamma\setminus\T\phi_0\bigr\}<1,
 \tag{3.1}\label{eq:gap}
\]
where $\sup\varnothing=0$.  The set $\Gamma$ is norming when
\[
 \norm y=\sup_{\gamma\in\Gamma}|\gamma(y)|
 \qquad (y\in Y).
\]
\end{definition}

\begin{theorem}\label{thm:main}
Let $X$ and $Y\ne\{0\}$ be Banach spaces, and suppose that $Y$ has a norming set $\Gamma$ with an isolated norming orbit witnessed by $(\phi_0,y_0)$.  Let $\mathcal I\subset\calL(X,Y)$ be a closed linear subspace such that
\[
 y_0\otimes f\in\mathcal I
 \qquad\text{for every }f\in X^*.
 \tag{3.2}\label{eq:rankone}
\]
If $\NA(X,Y)\cap\mathcal I$ is residual in $\mathcal I$, then $\NA(X,\K)$ is residual in $X^*$.
\end{theorem}

\begin{proof}
Set
\[
 \mathcal I_0=\{H\in\mathcal I:\phi_0\circ H=0\}.
\]
This is a closed subspace of $\mathcal I$.  By \eqref{eq:rankone}, the map
\[
 \Psi:X^*\times\mathcal I_0\longrightarrow\mathcal I,
 \qquad
 \Psi(f,H)=y_0\otimes f+H,
 \tag{3.3}\label{eq:Psi}
\]
is a linear homeomorphism.  Its inverse is
\[
 T\longmapsto
 \bigl(\phi_0\circ T,\,T-y_0\otimes(\phi_0\circ T)\bigr).
\]
Notice that every $H\in\mathcal I_0$ takes values in $Z:=\ker\phi_0$.

Put
\[
 C=\sup\bigl\{\norm{\gamma|_Z}:
       \gamma\in\Gamma\setminus\T\phi_0\bigr\}\le 1,
\]
again taking $C=0$ when the indexing set is empty.  Define
\[
 \mathcal W=\bigl\{\Psi(f,H):(1-\beta)\norm f>C\norm H\bigr\}.
 \tag{3.4}\label{eq:W}
\]
This is open.  Fix $T=\Psi(f,H)\in\mathcal W$.  If
$\gamma\in\Gamma\cap\T\phi_0$, then $\norm{\gamma\circ T}=\norm f$.  If
$\gamma\notin\T\phi_0$, then
\[
 \norm{\gamma\circ T}
 \le |\gamma(y_0)|\norm f+\norm{\gamma|_Z}\norm H
 \le \beta\norm f+C\norm H
 <\norm f.
 \tag{3.5}\label{eq:domination}
\]
Since $\Gamma$ norms $Y$,
\[
 \norm T
 =\sup_{\gamma\in\Gamma}\norm{\gamma\circ T}
 =\norm f.
 \tag{3.6}\label{eq:normequal}
\]
We claim that on $\mathcal W$,
\[
 T\in\NA(X,Y)
 \quad\Longleftrightarrow\quad
 f\in\NA(X,\K).
 \tag{3.7}\label{eq:equiv}
\]
If $|f(x)|=\norm f$ for some $x\in\Sphere_X$, then
\[
 \norm{Tx}\ge |\phi_0(Tx)|=|f(x)|=\norm T,
\]
so $T$ attains its norm.

Conversely, suppose that $\norm{Tx}=\norm T=\norm f$ for some
$x\in\Sphere_X$.  By \eqref{eq:domination},
\[
 \sup_{\gamma\in\Gamma\setminus\T\phi_0}|\gamma(Tx)|
 \le \beta\norm f+C\norm H<\norm f.
\]
For $\gamma\in\Gamma\cap\T\phi_0$, one has
$|\gamma(Tx)|=|f(x)|$.  Since
\[
 \norm{Tx}=\sup_{\gamma\in\Gamma}|\gamma(Tx)|=\norm f,
\]
it follows that $|f(x)|=\norm f$.  This proves \eqref{eq:equiv}.

Let $A=\NA(X,\K)$ and, for $m\ge1$, let
\[
 U_m=\{f\in X^*:\norm f>1/m\}.
\]
If $C>0$, set
\[
 r_m=\frac{1-\beta}{2Cm};
\]
if $C=0$, choose any $r_m>0$.  Then
\[
 U_m\times B_{\mathcal I_0}(0,r_m)
 \subset \Psi^{-1}(\mathcal W).
\]
Residuality of $\NA(X,Y)\cap\mathcal I$, restriction to this open cylinder, and \eqref{eq:equiv} show that
\[
 (A\cap U_m)\times B_{\mathcal I_0}(0,r_m)
\]
is residual in
\[
 U_m\times B_{\mathcal I_0}(0,r_m).
\]
By \cref{lem:cylinder}, $A\cap U_m$ is residual in $U_m$.  Hence
$U_m\setminus A$ is meagre in $X^*$ for every $m$.  Finally,
$0\in A$ and
\[
 X^*\setminus\{0\}=\bigcup_{m=1}^{\infty}U_m,
\]
so $X^*\setminus A$ is meagre.  Thus $A$ is residual in $X^*$.
\end{proof}

\begin{remark}
The theorem only needs one isolated orbit.  It does not require a finite boundary, separability of either space, or residuality in the whole operator space.  The role of the open cylinder is essential: restricting a residual set to the closed rank-one slice $\{y_0\otimes f:f\in X^*\}$ would be invalid in general.
\end{remark}

\section{Consequences}

\subsection{The proposed finite-boundary theorem}

\begin{corollary}\label{cor:finiteboundary}
Suppose there is a finite set $\Phi\subset\Sphere_{Y^*}$ such that
\[
 \norm y=\max_{\phi\in\Phi}|\phi(y)|
 \qquad(y\in Y),
\]
and that some $\phi_0\in\Phi$ admits $y_0\in\Sphere_Y$ with
\[
 \phi_0(y_0)=1,
 \qquad
 \max_{\phi\in\Phi\setminus\T\phi_0}|\phi(y_0)|<1.
\]
If $\NA(X,Y)$ is residual in $\calL(X,Y)$, then $\NA(X,\K)$ is residual in $X^*$.
\end{corollary}

\begin{proof}
Apply \cref{thm:main} with $\Gamma=\Phi$ and $\mathcal I=\calL(X,Y)$.
\end{proof}

\begin{corollary}\label{cor:polyhedral}
Let $Y$ be a nonzero finite-dimensional real polyhedral Banach space.  If $\NA(X,Y)$ is residual in $\calL(X,Y)$, then $\NA(X,\mathbb R)$ is residual in $X^*$.
\end{corollary}

\begin{proof}
Choose one representative from each pair of opposite vertices of $B_{Y^*}$.  These representatives form a finite norming set.  Every vertex of a finite-dimensional polytope is exposed.  If $\phi_0$ is exposed by $y_0$, normalize so that $\phi_0(y_0)=1$.  Exposure of both $\phi$ and $-\phi$ gives
$|\phi(y_0)|<1$ for every other orbit.  Now use \cref{cor:finiteboundary}.
\end{proof}

\subsection{Ranges with property \texorpdfstring{quasi-$\beta$}{quasi-beta}}

We recall the part of property quasi-$\beta$ that is needed.  There are a set
$\mathcal A\subset\Sphere_{Y^*}$, a map
$\sigma:\mathcal A\to\Sphere_Y$, and numbers $\rho(a)<1$ such that
\begin{align*}
 a(\sigma(a))&=1,\\
 |b(\sigma(a))|&\le\rho(a)\qquad(a,b\in\mathcal A,\ a\ne b),
\end{align*}
and every extreme point of $B_{Y^*}$, after multiplication by an element of $\T$, lies in the weak-star closure of a subset of $\mathcal A$; see \cite{AAP} or \cite[Section~3.1]{JMRZ}.

\begin{corollary}\label{cor:quasibeta}
Let $Y\ne\{0\}$ have property quasi-$\beta$.  If $\NA(X,Y)$ is residual in $\calL(X,Y)$, then $\NA(X,\K)$ is residual in $X^*$.
\end{corollary}

\begin{proof}
First, $\mathcal A$ is norming.  Indeed, for $y\in Y$ and
$e^*\in\ext(B_{Y^*})$, choose $t\in\T$ and a net $(a_i)$ in the corresponding subset of $\mathcal A$ such that $a_i\to te^*$ in the weak-star topology.  Then
\[
 |e^*(y)|=\lim_i|a_i(y)|
 \le\sup_{a\in\mathcal A}|a(y)|.
\]
Taking the supremum over the extreme points and using Krein--Milman gives
$\norm y\le\sup_{a\in\mathcal A}|a(y)|$; the reverse inequality is immediate.

Fix $a_0\in\mathcal A$ and put $y_0=\sigma(a_0)$.  No distinct rotation of $a_0$ can belong to $\mathcal A$, since such an element would have modulus one at $y_0$, contradicting the second quasi-$\beta$ inequality.  Hence
\[
 \sup_{a\in\mathcal A\setminus\T a_0}|a(y_0)|
 \le\rho(a_0)<1.
\]
Thus $\mathcal A$ has an isolated norming orbit, and \cref{thm:main} applies.
\end{proof}

In particular, \cref{cor:quasibeta} covers property $\beta$, all finite-dimensional polyhedral ranges, real polyhedral preduals of $\ell_1$, and arbitrary real or complex closed subspaces of $c_0(\Gamma)$; the latter examples are recorded in \cite[Example~3.2]{JMRZ}.

\subsection{Finite-dimensional ranges with countably many extreme orbits}

The next consequence is strictly wider than the finite-boundary case.

\begin{theorem}\label{thm:countableorbits}
Let $Y\ne\{0\}$ be finite-dimensional and suppose that
\[
 \ext(B_{Y^*})/\T
\]
is countable.  Then $Y$ has a norming set with an isolated norming orbit.  Consequently, for every Banach space $X$,
\[
 \NA(X,Y)\text{ residual in }\calL(X,Y)
 \quad\Longrightarrow\quad
 \NA(X,\K)\text{ residual in }X^*.
\]
\end{theorem}

\begin{proof}
Put $K=B_{Y^*}$ and $E=\ext(K)$.  We regard the finite-dimensional space $Y^*$ as a real normed space when discussing convexity.  The set $E$ is a $G_\delta$ subset of $K$.  Indeed,
\[
 K\setminus E=\bigcup_{n=1}^{\infty}F_n,
\]
where
\[
 F_n=\left\{z\in K:\ z=\frac{u+v}{2}
       \text{ for some }u,v\in K\text{ with }\norm{u-v}\ge\frac1n\right\}.
\]
Each $F_n$ is compact.  Hence $E$ is completely metrizable and therefore a Baire space.

By hypothesis, $E$ is the union of countably many $\T$-orbits.  Each orbit is compact and thus closed in $E$.  If every orbit had empty interior in $E$, then $E$ would be meagre in itself, a contradiction.  Therefore one orbit
$\mathcal O=\T\phi_0$ is open in $E$.

Let
\[
 F=\overline{E\setminus\mathcal O}^{\,K}.
\]
The openness of $\mathcal O$ in $E$ implies $F\cap\mathcal O=\varnothing$.  If $E\setminus\mathcal O=\varnothing$, choose a norm-one vector $y_0$ at which $\phi_0$ attains its norm; the conclusion is immediate.  Otherwise let
\[
 C=\co(F).
\]
This is compact, convex, and invariant under multiplication by $\T$.  Moreover,
$\phi_0\notin C$: if an extreme point $\phi_0$ were a convex combination of points of $F\subset K$, every point occurring with positive coefficient would have to equal $\phi_0$, contradicting $F\cap\mathcal O=\varnothing$.

Strict separation in the underlying real space gives $y\in Y$ such that
\[
 \operatorname{Re}\phi_0(y)>
 \sup_{\psi\in C}\operatorname{Re}\psi(y).
\]
In the complex case, every real-linear functional on $Y^*$ has the form
$\psi\mapsto\operatorname{Re}\psi(y)$ because $Y$ is finite-dimensional.  Since $C$ is $\T$-invariant,
\[
 s:=\sup_{\psi\in C}\operatorname{Re}\psi(y)
   =\sup_{\psi\in C}|\psi(y)|.
\]
Rotate $\phi_0$ inside its orbit so that
$a:=\phi_0(y)=|\phi_0(y)|>s$.  As $E$ is norming and every extreme point belongs either to $\mathcal O$ or to $F\subset C$,
\[
 \norm y=a.
\]
Set $y_0=y/a$.  Then
\[
 \phi_0(y_0)=1,
 \qquad
 \sup_{\psi\in E\setminus\mathcal O}|\psi(y_0)|\le s/a<1.
\]
Choose one representative from each $\T$-orbit in $E$, including $\phi_0$, and call the resulting set $\Gamma$.  Krein--Milman gives
\[
 \norm z=\sup_{\gamma\in\Gamma}|\gamma(z)|
 \qquad(z\in Y),
\]
so $\Gamma$ is norming and $(\phi_0,y_0)$ witnesses an isolated norming orbit.  The final implication is \cref{thm:main}.
\end{proof}

Combining \cref{thm:countableorbits} with \cite[Example~3.18(a)]{JMRZ} yields the following equivalence.

\begin{corollary}\label{cor:equivalence}
Let $Y\ne\{0\}$ be finite-dimensional with
$\ext(B_{Y^*})/\T$ countable.  For every Banach space $X$,
\[
 \NA(X,Y)\text{ is residual in }\calL(X,Y)
 \quad\Longleftrightarrow\quad
 \NA(X,\K)\text{ is residual in }X^*.
\]
\end{corollary}

\section{Why this does not settle the full question}

The proof uses a genuine geometric gap.  If $Y$ is smooth and its norming functionals vary continuously, there need not be any isolated norming orbit.  The real Euclidean plane is the basic example: its dual extreme points form a circle, and no orbit modulo signs is isolated.  In that setting an arbitrarily small transverse perturbation can rotate the norming functional, so there is no open cylinder on which vector norm attainment is exactly the same as norm attainment of one fixed scalar coordinate.

Several natural routes stop at this point.
\begin{enumerate}[label=(\roman*),leftmargin=2.2em]
\item Restricting a residual subset of $\calL(X,Y)$ to a rank-one subspace is invalid because that subspace is closed and typically nowhere open.
\item Perturbing $y_0\otimes f$ by operators tending to zero produces norm-attaining scalar combinations converging to $f$, but this gives no new information: $\NA(X,\K)$ is always dense by Bishop--Phelps and need not be norm closed.
\item A compactness argument may produce norming functionals $\phi_n\to\phi$, but the corresponding scalar norm-attainers $T^*\phi_n$ can converge to a non-norm-attaining functional.  This is precisely where the uniform gap in \eqref{eq:gap} is used.
\item For a smooth finite-dimensional range there is an uncountable continuum of possible norming directions, so a countable Baire-category pigeonhole argument does not select a fixed coordinate.
\end{enumerate}

The recent construction of a Banach space whose norm-attaining functionals are algebraically very small \cite{Martin} suggests possible counterexample domains, but it does not currently supply a counterexample here.  In fact, closely related existence questions for rank-two norm-attaining operators into the two-dimensional Hilbert space remain delicate \cite{Martin,KLMW}.

\section{Conclusion}

The submitted partial theorem is mathematically sound after one repair: replace the unjustified general Kuratowski--Ulam invocation by the cylinder-reflection theorem for completely metrizable spaces, and explicitly record the equality $\norm T=\norm{\phi_0\circ T}$ in the domination region.

The same mechanism proves substantially more.  Residual vector-valued norm attainment forces residual scalar norm attainment whenever the range has one isolated norming orbit.  This includes property quasi-$\beta$ and all finite-dimensional ranges with countably many extreme dual orbits.  No complete proof or counterexample for arbitrary nonzero $Y$ was located, and smooth ranges such as the Euclidean plane remain outside the method.

\begin{thebibliography}{99}

\bibitem{AAP}
M.~D. Acosta, F.~J. Aguirre and R.~Pay\'{a},
\emph{A new sufficient condition for the denseness of norm attaining operators},
Rocky Mountain J. Math. \textbf{26} (1996), 407--418.

\bibitem{Brian}
W. Brian,
\emph{Cardinal invariants of a meager ideal},
Topology Appl. \textbf{379} (2026), Article 109516.

\bibitem{RSE}
G. Choi, H. del R\'{\i}o, A. Fovelle, M. Jung and M. Mart\'{\i}n,
\emph{Range strongly exposing operators between Banach spaces},
arXiv:2503.18581v2 (2025).

\bibitem{GMZ}
A.~J. Guirao, V. Montesinos and V. Zizler,
\emph{Remarks on the set of norm-attaining functionals and differentiability},
Studia Math. \textbf{241} (2018), 71--86.

\bibitem{JMRZ}
M. Jung, M. Mart\'{\i}n and A. Rueda Zoca,
\emph{Residuality in the set of norm attaining operators between Banach spaces},
J. Funct. Anal. \textbf{284} (2023), Article 109746.

\bibitem{KLMW}
V. Kadets, G. L\'{o}pez, M. Mart\'{\i}n and D. Werner,
\emph{Norm attaining operators of finite rank},
in \emph{The Mathematical Legacy of Victor Lomonosov},
Advances in Analysis and Geometry, vol.~2, de Gruyter, 2020, 157--187.

\bibitem{Martin}
M. Mart\'{\i}n,
\emph{A Banach space whose set of norm-attaining functionals is algebraically trivial},
J. Funct. Anal. \textbf{288} (2025), Article 110815.

\bibitem{MT}
W.~B. Moors and N.~O. Tan,
\emph{Dual differentiation spaces},
Bull. Aust. Math. Soc. \textbf{99} (2019), 467--472.

\end{thebibliography}

\end{document}
